\newpage
\vspace*{7cm}
\begin{center}
  {\Huge\bfseries CONCLUSION GÉNÉRALE}
\end{center}
\vspace*{\fill}
\addcontentsline{toc}{chapter}{CONCLUSION GÉNÉRALE}
\markboth{CONCLUSION GÉNÉRALE}{CONCLUSION GÉNÉRALE}

\newpage

Ce projet académique et expérimental avait pour vocation première de jauger, sous un prisme strictement qualitatif et quantitatif orienté vers la clinique, les performances respectives de trois piliers fondamentaux de l'apprentissage automatique supervisé (la modélisation par Vecteurs de Support SVM, la construction ensembliste des Forêts Aléatoires et l'optimisation par Descente de Gradient Extrême XGBoost) dans le cadre redoutable de la prédiction systémique précoce du diabète de type 2. Cette analyse a pris pour fondation exclusive la base de données de référence Pima Indians Diabetes Database.

% ---- Synthèse des contributions ----
\section*{Synthèse des contributions et approfondissements théoriques}

La \textbf{Partie I} de l'ouvrage s'est attelée à bâtir l'infrastructure théorique requise pour saisir les subtilités mécaniques présidant aux déductions des algorithmes choisis. Au sujet des SVM, nous avons déployé rigoureusement le corpus mathématique gérant les hyperplans à marge maximale, conjugué à la résolution du problème dual de Lagrange et à l'astuce de projection spatiale (noyau), en expliquant par la théorie les défaillances inhérentes de la méthode face aux densités tabulaires limitées. Concernant les Forêts Aléatoires, l'examen s'est focalisé sur la germination des arbres de classification, magnifiée par les lois de la sélection stochastique garantissant une inestimable robustesse face au surapprentissage. Pour finir, le spectre de l'algorithme XGBoost a été exploré en décortiquant minutieusement l'approche d'approximation du second ordre de la fonction de perte (la formule de Taylor) qui, lorsqu'elle est assortie à la répression structurelle des nœuds, garantit un apprentissage virtuellement inattaquable.

L'étude liminaire (chapitre 4) a de surcroît mis en évidence la cacophonie méthodologique régnant au sein des publications antérieures, dénonçant avec vigueur l'hégémonie de métriques trompeuses (comme la précision globale prise isolément) et les graves insuffisances touchant tant à l'explicabilité du modèle qu'à sa propension à se généraliser en dehors de son contexte initial.

La \textbf{Partie II} a, pour sa part, matérialisé un protocole expérimental à la fois souverain et intolérant aux approximations, s'illustrant par les accomplissements suivants :

\begin{enumerate}
    \item \textbf{Un assainissement chirurgical de la matrice d'origine}, recourant à une substitution par la médiane segmentée selon la classe cible, secondée par la neutralisation des valeurs aberrantes aux percentiles asymétriques extrêmes, ainsi qu'une normalisation mathématique profonde (Yeo-Johnson) imposée par la nature distordue des données médicales.
    
    \item \textbf{Une ingénierie biomathématique de très haut vol}, donnant naissance à 8 indicateurs synthétiques greffés sur le socle originel. L'incidence du paramètre \texttt{Risk\_Score} (corrélation de Pearson atteignant $+0.535$) et de la collusion \texttt{Glucose\_BMI} ($+0.526$) certifie la nécessité absolue de sceller l'alliance entre pragmatisme médical et modélisation algorithmique, portant l'architecture finale à 16 facteurs d'intérêts.
    
    \item \textbf{Une riposte structurelle au déficit des diagnostics pathologiques certifiés}, s'incarnant par l'implémentation du protocole Borderline-SMOTE, permettant de symétriser numériquement les profils simulés (400 face à 400 dans le lot d'apprentissage) sans perturber fondamentalement la topologie de l'espace de décision.
    
    \item \textbf{Une traversée d'optimisation probabiliste orchestrée par Optuna} (englobant 200 évaluations massives concernant XGBoost), traçant la voie vers des configurations paramétriques d'exception hors de portée d'une simple exploration naïve.
\end{enumerate}

% ---- Bilan des résultats ----
\section*{Bilan chiffré des investigations}

Les expérimentations attestent catégoriquement de l'effondrement des architectures isolées au bénéfice absolu des procédés de modélisation ensembiliste. Au faîte de la hiérarchie technique, \textbf{XGBoost} et \textbf{Random Forest} coiffent conjointement les sommets prédictifs avec une exactitude figée à 89,0\% associée à un indice F1 de 0.844. Sur le versant évaluant l'homogénéité de discrimination, ils brisent les plafonds avec des aires colossales respectives fixées à \textbf{0.9496} et \textbf{0.9453}. L'écart béant instauré par rapport aux manipulations sans relief observées ordinairement (+20,9 points d'ascension sur l'indice F1-Score comparativement à l'état originel des données) entérine solennellement le succès de l'expérimentation.

Malgré des paramétrages poussés dans ses retranchements optimaux, l'hyperplan vectoriel (SVM) ne hissera son bouclier qu'à 0.859 d'aire discriminatoire, accablé par un index de rappel dramatiquement insuffisant. Cette asphyxie mécanique prouve son inadaptation foncière pour s'accommoder des intrications masquées nichées dans des tableaux cliniques moyennement denses.

En outre, il est établi que la surabondance systémique (méta-modèles ou coalitions de votes) s'est avérée stérile pour extirper plus de sens que le parangon de la session (XGBoost), bloquée par l'inévitable corrélation partagée par les pronostics des modèles de premier plan.

% ---- Limites et perspectives ----
\section*{Carences intrinsèques et desseins futurs}

Si le triomphe de l'expérimentation impressionne, il convient de ne pas se méprendre sur ses assises périlleuses. La déduction algorithmique restreinte de manière claustrale aux seules amérindiennes Pimas ne saurait octroyer de sauf-conduit universel. D'autre part, la reconstruction spéculative frôlant les 48,7\% de l'attribut majeur (Insuline) instaure indéniablement un artéfact statistique lourd de conséquences potentielles. Enfin, le silence opaque et imperturbable entourant les ramifications décisionnelles intimes de l'algorithme (manque d'explicabilité par profil individuel) reste pénalisant d'un point de vue médicolégal.

Surgissent alors des futurs chantiers de recherche majeurs orientés vers l'adhésion d'architectures tabulaires d'explicabilité (à l'instar des protocoles mathématiques préconisés par Shapley) pour légitimer la sentence prédictive aux yeux du praticien humain. Ils appelleront de surcroît à sonder un pan entier de biométries non intrusives, esquissant l'horizon d'un pronostic de masse indolore et accessible techniquement à l'échelle transcontinentale.

% ---- Mot final ----
\section*{Dernière observation algorithmique}

L'investigation vient de sceller sans équivoque un postulat majeur de la science de la donnée contemporaine. L'élection fastueuse d'un solveur de dernière génération importe intrinsèquement beaucoup moins que l'acuité et le ciselage investis à métamorphoser la substance primitive en vecteurs clairs, balancés et cliniquement chargés d'instructions. Au carrefour exact de la médecine interne et des statistiques matricielles rigoureuses, cet ouvrage a démontré qu'une intelligence prédictive infaillible exige avant tout des données rigoureusement respectées et modélisées avec la plus haute méticulosité scientifique.